\documentclass[11pt]{article}
\usepackage[a4paper,margin=24mm]{geometry}
\usepackage[T1]{fontenc}
\usepackage{lmodern,amsmath,amssymb,amsthm,mathtools,microtype}
\usepackage[hidelinks]{hyperref}
\usepackage{xurl}
\newtheorem{theorem}{Theorem}
\newtheorem{lemma}{Lemma}
\newtheorem{remark}{Remark}
\newcommand{\R}{\mathbb R}
\newcommand{\dd}{\,\mathrm d}
\newcommand{\norm}[1]{\lVert#1\rVert}
\newcommand{\ip}[2]{\langle#1,#2\rangle}
\newcommand{\ind}{\mathbf1}
\newcommand{\HH}{\mathcal H}
\newcommand{\RR}{\mathcal R}
\setlength{\emergencystretch}{2em}
\allowdisplaybreaks[1]
\title{A finite-rank Schur correction from a weighted prime defect}
\author{Jongmin Choi\\Independent Researcher, Seoul, Korea\\
ORCID iD: \href{https://orcid.org/0009-0008-7448-514X}{0009-0008-7448-514X}}
\date{28 October 2026\quad Analytic research note v0.1}
\hypersetup{pdftitle={A finite-rank Schur correction from a weighted prime defect},pdfauthor={Jongmin Choi}}
\begin{document}
\maketitle
\begin{abstract}
At support radius $B=3/4$, we retain a positive remainder of the complete prime operator with active prime powers $2,3,4$. A bounded weighted edge identity gives $E=RI-P\succeq I/54$, with an explicit $R<27/20$. A localized observation of this remainder is converted into an exact finite cosine feature by paying for its omitted tail in the same energy norm. This yields a rank-one improvement of the diagonal-inverse Schur bound: the coupling cost is at most its full baseline bound minus $(8/25)\zeta\zeta^*$, where $\zeta$ is an explicit finite sum. The baseline infinite coupling tail remains present. No certified nonvanishing of the selected vector or positivity of the corrected low matrix is reported. A separate fixed-parameter implementation failed its positivity gate. The result is a fixed-radius analytic inequality, not a new local positivity certificate, an all-radius theorem, or a proof of the Riemann hypothesis.
\end{abstract}

\section{The signed form and the operator split}
All norms are ordinary $L^2$ norms. Let $\Lambda$ be the von Mangoldt function and set
\[
 \rho(v)=\frac{e^{-v/2}}{1-e^{-2v}},\quad a_n=\frac{\Lambda(n)}{\sqrt n},\quad
 c_\Gamma=\log\pi+\gamma_E+\frac\pi2+3\log2.
\]
For real even smooth functions supported in $(-B,B)$, extended by zero, define
\begin{align}
 t_\Gamma(u)&=\int_0^\infty\rho(v)\norm{u(\cdot+v)-u}_2^2\dd v,\label{eq:gamma}\\
 q_B(u)&=t_\Gamma(u)-c_\Gamma\norm u_2^2
 -2\sum_{n\ge2}a_n\int_\R u(x)u(x+\log n)\dd x\nonumber\\
 &\hspace{25mm}+2\left|\int_\R\cosh(x/2)u(x)\dd x\right|^2.\label{eq:form}
\end{align}
The minimum local domain $\mathcal D_B$ is the completion of these tests for the norm $(t_\Gamma+\norm{\cdot}_2^2)^{1/2}$. On fixed support, the prime correlations are a finite sum of bounded forms and the pole is bounded. Thus \eqref{eq:form} extends continuously to $\mathcal D_B$.

Fix $B=3/4$ throughout. The unitary identification $f(x)=\sqrt2u(x)$ for $0<x<B$ puts the form on $\HH=L^2(0,B)$. Set
\[
 s=\log2,\quad t=\log3,\quad d=B-s,\qquad
 a=\frac{s}{\sqrt2},\quad b=\frac{s}{2},\quad c=\frac{t}{\sqrt3}.
\]
Exactly $2,3,4$ are active, since $\log4<2B<\log5$; all later correlations are zero by support. In particular the weight of $4$ is $b=\log2/2$, not $\log4/2$.

The associated bounded self-adjoint prime operator $P=P_2+P_3+P_4$ is
\begin{align}
 P_2f(x)&=a\{\ind_{(0,d)}(x)f(x+s)+\ind_{(s,B)}(x)f(x-s)
                  +\ind_{(0,s)}(x)f(s-x)\},\label{eq:prime}\\
 P_3f(x)&=c\ind_{(t-B,B)}(x)f(t-x),\nonumber\\
 P_4f(x)&=b\ind_{(2s-B,B)}(x)f(2s-x).\nonumber
\end{align}
These formulas retain both same-side translations and the reflection of prime $2$. Their quadratic forms equal the complete prime term in \eqref{eq:form} under the unitary identification.

\section{The weighted prime remainder}
Partition $(0,B)$, up to endpoints, into
\[
 C=(0,d),\quad L_-=(s-d,s),\quad L_+=(s,s+d),\quad
 L=L_-\cup L_+,\quad J=(d,s-d).
\]
Prime $2$ joins each $y\in C$ to $s-y$ and $s+y$, and reflects $J$ about $s/2$. Prime $4$ joins the two leaves. Put $S=(t-B,B)$. Every leaf lies in $S$, and
\begin{equation}
 t-L=(t-B,t-2s+B)\subset J.\label{eq:geometry}
\end{equation}
Indeed $2s>t$, $s+t>2B$, and $t+2B<4s$. The last two follow already from $s>2/3$, $1<t<11/10$. This leaf-to-middle property is specific to the present radius.

Define
\begin{equation}
 R=\frac{b+c+\sqrt{(b+c)^2+8a^2}}2,\qquad
 \eta=\frac{2a}{R},\qquad w=\eta\ind_C+\ind_{J\cup L},\qquad E=RI-P.
 \label{eq:weight}
\end{equation}
The polynomial $F(x)=x^2-(b+c)x-2a^2$ vanishes at $R$. The elementary bounds
\begin{equation}
 a>\frac7{15},\quad a<\frac12,\quad
 \frac47<c<\frac{121}{190},\quad \frac54<R<\frac{27}{20},\quad 0<\eta<1
 \label{eq:rational}
\end{equation}
will be useful. Here $s<7/10$, $t<11/10$ and $\sqrt3>19/11$ imply $b+c<75/76$ and $2a^2<49/100$, while
\[
 F(5/4)<\frac{25}{16}-\frac{95}{84}-\frac49=-\frac{13}{1008},
 \qquad
 \left(\frac{27}{20}\right)^2-\frac{75}{76}\frac{27}{20}-\frac{49}{100}
 =\frac1{3800}.
\]
Moreover $b+c>a$ implies $R>2a$, and
\begin{equation}
 \delta:=R-a-c>\frac54-\frac12-\frac{121}{190}=\frac{43}{380}.
 \label{eq:delta}
\end{equation}
The logarithmic bounds follow from the positive series for $\log((1+x)/(1-x))$ and its geometric tail; in particular $s>56/81>2/3$.

\begin{lemma}[Exact bounded edge identity]\label{lem:edge}
For every real $f\in\HH$,
\begin{align}
 \ip f{Ef}={}&a\eta\int_0^d\left(
 \left|\frac{f(y)}\eta-f(s-y)\right|^2+
 \left|\frac{f(y)}\eta-f(s+y)\right|^2\right)\dd y\nonumber\\
 &+b\int_0^d|f(s-y)-f(s+y)|^2\dd y
 +\frac a2\int_J|f(x)-f(s-x)|^2\dd x\nonumber\\
 &+\frac c2\int_S|f(x)-f(t-x)|^2\dd x
 +\int_J\left(R-a-c\ind_S(x)\right)|f(x)|^2\dd x.\label{eq:edge}
\end{align}
Consequently $0\preceq E\preceq2RI$ and $\norm P\le R$.
\end{lemma}
\begin{proof}
The weighted row ratio $(Pw)/w$ equals $2a/\eta=R$ on $C$, $a\eta+b+c=R$ on $L$, and $a+c\ind_S$ on $J$. An unoriented edge of weight $p$ contributes
$pw(x)w(y)|f(x)/w(x)-f(y)/w(y)|^2$.
Expansion gives the two weighted diagonal terms and $-2pf(x)f(y)$. The factors $1/2$ in the two reflection integrals avoid counting an edge twice. The remaining diagonal is precisely $R-(Pw)/w$. This proves the identity first on bounded functions and then by $L^2$ continuity. Applying the same edge estimate to $|f|$ bounds the absolute value of $\ip f{Pf}$ by $R\norm f^2$, giving both operator inequalities.
\end{proof}

The identity is a bounded graph version of a ground-state remainder representation; compare \cite[Proposition 2.3 and Remark 2.4]{FS}. We proved it directly because the prime shifts are atomic, not a kernel density satisfying a continuous-kernel hypothesis. In particular, we never apply the unbounded Gamma form to $f/w$.

\begin{theorem}[Coercivity of the prime remainder]\label{thm:gap}
At $B=3/4$, the bounded operator in \eqref{eq:weight} satisfies $E\succeq I/54$.
\end{theorem}
\begin{proof}
Write $e=\ip f{Ef}$ and let $C_f,L_f,J_f$ be the squared norms on $C,L,J$. The middle killing in \eqref{eq:edge} gives $J_f\le e/\delta$. By \eqref{eq:geometry}, the leaf-to-middle matching is injective and measure preserving, so
\[
 L_f\le\frac{2e}{c}+2J_f.
\]
Averaging the two center-to-leaf square inequalities gives
\[
 C_f\le\frac\eta a e+\eta^2L_f.
\]
Hence, for nonzero $f$,
\[
 \norm f^2\le\left(\frac\eta a+\frac{2(1+\eta^2)}c+
                    \frac{3+2\eta^2}\delta\right)e
 <\left(\frac{15}{7}+7+\frac{1900}{43}\right)e
 =\frac{16052}{301}e<54e.
\]
The zero function is immediate. Each retained nonnegative energy was bounded above by the same total $e$; the displayed combined coefficient makes this conservative accounting explicit.
\end{proof}

\section{A finite observation of the positive remainder}
Use the orthonormal cosine basis
\[
 e_n(x)=b_n\cos(k_nx),\qquad b_n=\sqrt{\frac{2-\delta_{n0}}B},
 \qquad k_n=\frac{\pi n}{B}.
\]
Let $\Pi_H$ project onto modes $n\ge M=32$ and put $N=16384$. The interval $I=[1/8,1/4]$ lies inside $(d,t-B)\subset J$. There the last integrand in \eqref{eq:edge} has coefficient $R-a>c>1/2$. Thus, with $\psi=\sqrt8\ind_I$,
\begin{equation}
 \norm\psi=1,\qquad E\succeq\tfrac12|\psi\rangle\langle\psi|.
 \label{eq:observation}
\end{equation}
Here the rank-one notation means $\ip f{|\psi\rangle\langle\psi|f}=|\ip\psi f|^2$.
Define the exact finite feature
\begin{equation}
 \phi=\sum_{n=M}^{N-1}\psi_ne_n,\qquad
 \psi_n=\frac{\sqrt8 b_n}{k_n}\left[\sin(k_n/4)-\sin(k_n/8)\right].
 \label{eq:feature}
\end{equation}
No membership of the indicator $\psi$ in the Gamma form domain is needed.

\begin{lemma}[Finite feature with its tail paid for]\label{lem:feature}
On $H=\Pi_H\HH$, one has $E_H:=\Pi_HE\Pi_H\succeq(2/5)|\phi\rangle\langle\phi|$.
\end{lemma}
\begin{proof}
For $n\ge1$, $|\psi_n|^2\le48/(\pi^2n^2)<16/(3n^2)$. Therefore
\[
 \tau_N:=\norm{\Pi_{[N,\infty)}\psi}^2<\frac{16}{3(N-1)}.
\]
For $h\in H$, \eqref{eq:observation} and Theorem~\ref{thm:gap} give a single triangle inequality in the dual $E$-energy norm:
\[
 |\ip\phi h|\le |\ip\psi h|+|\ip{\Pi_{[N,\infty)}\psi}h|
 \le\left(\sqrt2+\sqrt{54\tau_N}\right)\ip h{Eh}^{1/2}.
\]
At $N=16384$,
\[
 54\tau_N<\frac{288}{16383}<\frac4{225},\qquad
 \left(\frac{10}{7}+\frac2{15}\right)^2
 =\frac{26896}{11025}<\frac52,
\]
since $\sqrt2<10/7$. Squaring proves the assertion. This argument derives one lower bound from two estimates of the same energy; it does not add two copies of that energy.
\end{proof}

\section{The finite Schur correction}
Split the exact form into its first $M$ cosine modes and the high space:
\begin{equation}
 q(f)=z^*Az+2\ip h{\RR z}+t_H[h],\qquad f=z+h.\label{eq:blocks}
\end{equation}
The real symmetric low matrix is $A$, the high operator is $T$, and the exact coupling has entries $r_{nm}=Q_{nm}$ for $n\ge M$, $m<M$. Appendix~\ref{app:gamma} derives, for the form \eqref{eq:form},
\begin{equation}
 T\succeq K+E_H,\qquad
 K=\operatorname{diag}_{n\ge M}(D_n-R)\succeq\frac85I.\label{eq:high}
\end{equation}
The positive Gamma boundary and pole high blocks are omitted only once in this inequality. The actual low and cross blocks are unchanged.

\begin{theorem}[Explicit finite subtraction]\label{thm:schur}
Define the full baseline coupling cost and a finite vector by
\begin{equation}
 C_0=\RR^*K^{-1}\RR,\qquad
 \zeta_m=\sum_{n=M}^{N-1}\frac{\psi_n r_{nm}}{D_n-R}.
 \label{eq:zeta}
\end{equation}
For every proved matrix majorant $\mathcal C\succeq C_0$,
\begin{equation}
 \RR^*T^{-1}\RR\preceq C_0-\frac8{25}\zeta\zeta^*
 \preceq\mathcal C-\frac8{25}\zeta\zeta^*.
 \label{eq:main}
\end{equation}
Consequently the exact low Schur complement is bounded below by
$A-\mathcal C+(8/25)\zeta\zeta^*$.
\end{theorem}
\begin{proof}
Lemma~\ref{lem:feature} gives $T\succeq K+(2/5)|\phi\rangle\langle\phi|$. Inverse monotonicity and the rank-one inverse formula yield
\[
 \RR^*T^{-1}\RR\preceq C_0-
 \frac{\zeta\zeta^*}{5/2+\chi},\qquad
 \chi=\ip\phi{K^{-1}\phi},\quad \zeta=\RR^*K^{-1}\phi.
\]
Since $\norm\phi\le1$ and $K\succeq(8/5)I$, one has $0\le\chi\le5/8$. The reciprocal denominator is therefore at least $8/25$. Both $\chi$ and $\zeta$ are exact finite sums. The final comparison follows by adding $\mathcal C-C_0\succeq0$.
\end{proof}

This is a concrete low-rank inverse comparison within the Schur reduction; compare the general framework in \cite[Theorem 1.2]{DSS}. The cited theorem does not assert positivity of this particular prime--Gamma form. The finite vector in \eqref{eq:zeta} has no omitted overlap tail: its feature tail was already accounted for in Lemma~\ref{lem:feature}. In contrast, $C_0$ is still an infinite high-mode sum. Replacing it by a finite row sum alone is not permitted. Any numerical majorant $\mathcal C$ must include every omitted high row.

The correction is nonzero exactly when $\zeta\ne0$. A certified interval evaluation can verify this by enclosing one component away from zero. Neither that nonvanishing nor positive definiteness of $A-\mathcal C+(8/25)\zeta\zeta^*$ is asserted here. Entrywise interval bounds are not Loewner inequalities: one must enclose the entries of the exact corrected matrix, including signed products, and then apply a valid matrix positivity certificate.

\section{Strictness and resource accounting}
The entire remainder provides a useful alternative. From $0\preceq E_H\preceq(27/10)I$ and \eqref{eq:high}, the bounded operator $X=K^{-1/2}E_HK^{-1/2}$ satisfies $0\preceq X\preceq27/16$. Functional calculus gives
\begin{equation}
 T^{-1}\preceq K^{-1}-\frac{16}{43}K^{-1}E_HK^{-1}.
 \label{eq:inverse}
\end{equation}
Indeed $(1+x)^{-1}\le1-(16/43)x$ on $[0,27/16]$. No commutation of $K$ and $E_H$ is assumed. The resulting cost reduction is at least
\[
 \frac8{1161}\RR^*K^{-2}\RR,
\]
using $E_H\succeq I/54$. It is strictly positive in any low direction with nonzero actual high coupling. Such a direction exists: the Gamma image of the zero-extended constant mode has the endpoint term
\[
 b_0\left[\int_{B-x}^\infty\rho(v)\dd v+
                \int_{B+x}^\infty\rho(v)\dd v\right],
\]
which diverges logarithmically as $x\uparrow B$. The mass, prime and pole images are bounded and cannot cancel it. No finite cosine combination removes this logarithm, so $\Pi_HQe_0\ne0$. This does not imply that its overlap with the particular finite feature $\phi$ is nonzero.

Equations~\eqref{eq:main} and \eqref{eq:inverse} are alternative uses of the same remainder. Their improvements cannot simply be added. Similarly, the bounds $E\succeq I/54$ and $E_H\succeq(2/5)|\phi\rangle\langle\phi|$ cannot be added without a separate allocation proof. The dual-norm step in Lemma~\ref{lem:feature} is legitimate because it establishes just one new bound.

The geometry, observation interval and constants in this note are fixed-radius statements. Larger radii activate additional prime powers and change the graph. No all-radius conclusion follows. No historical novelty, human peer review or proof-assistant certification is claimed. The accompanying rational-arithmetic checks verify only the listed scalar manipulations, not the analytic proof or any numerical Schur certificate.

\appendix
\section{Gamma decomposition and the application hypotheses}\label{app:gamma}
For $\alpha_j=2j+1/2$, expand $\rho(v)=\sum_{j\ge0}e^{-\alpha_jv}$. Let $A_N$ be the Neumann Laplacian on $(0,B)$. Its Green kernel satisfies, for $\alpha>0$,
\begin{align*}
 e^{-\alpha|x-y|}+e^{-\alpha(x+y)}
 ={}&2\alpha(A_N+\alpha^2)^{-1}(x,y)\\
 &-2(\coth(\alpha B)-1)\cosh(\alpha x)\cosh(\alpha y).
\end{align*}
This follows by solving the elementary one-dimensional Neumann Green equation. Integrating one exponential shift form gives $2\norm f^2/\alpha$ minus the quadratic form of the kernel on the left. Thus, by summing its nonnegative decomposition,
\begin{align}
 t_\Gamma(u)&=\sum_{n\ge0}d_\Gamma(k_n)|f_n|^2+\mathcal B[f],\nonumber\\
 d_\Gamma(k)&=\sum_{j\ge0}\frac{2k^2}{\alpha_j(\alpha_j^2+k^2)},\qquad
 \mathcal B[f]=\sum_{j\ge0}2(\coth(\alpha_jB)-1)
       |\ip f{\cosh(\alpha_jx)}|^2\ge0.\label{eq:decomp}
\end{align}
With $D_n=d_\Gamma(k_n)-c_\Gamma$, the exact form is diagonal $D$, plus $\mathcal B$, plus the positive pole, minus $P$. Hence its high block obeys \eqref{eq:high} once the scalar floor is proved.

For completeness, the following source-free estimate suffices:
\begin{equation}
 d_\Gamma(k)-c_\Gamma\ge\log\frac{k}{2\pi}-\frac2{25}-\frac7{2k^2},\qquad k>0.
 \label{eq:symbol}
\end{equation}
Write $C_\Gamma=c_\Gamma-\log\pi$. The elementary harmonic identity
\[
 C_\Gamma=\lim_{L\to\infty}\left(\sum_{j=0}^{L-1}\frac1{j+1/4}-\log L\right)
\]
follows by splitting the harmonic sum into residue classes and using the alternating arctangent series. The trapezoidal remainder for $1/x$ on $[5/4,\infty)$ gives
\[
 0\le\varepsilon_0:=C_\Gamma-4+\log(5/4)-2/5\le2/25.
\]
In fact the sum of interval errors is bounded by $1/(8(5/4)^2)$. The positive convex function
$F_k(\alpha)=2k^2/[\alpha(\alpha^2+k^2)]$ then gives, by the mesh-two trapezoidal bound,
\[
 d_\Gamma(k)\ge F_k(1/2)+\tfrac12\log(1+k^2/(5/2)^2)+\tfrac12F_k(5/2).
\]
Substitution of $\varepsilon_0$ makes the cancellation explicit:
\begin{align*}
 d_\Gamma(k)-c_\Gamma\ge{}&\log\frac{k}{2\pi}-\varepsilon_0
 +\frac12\log\left(1+\frac{25}{4k^2}\right)\\
 &-\frac1{k^2+1/4}-\frac{5/2}{k^2+25/4}.
\end{align*}
Discarding the nonnegative logarithm and bounding the last two terms proves \eqref{eq:symbol}.

At $M=32$, $k_M=128\pi/3>128$, and
\[
 \log\frac{k_M}{2\pi}=6\log2-\log3>\frac{112}{27}-\frac{11}{10}.
\]
Together with $R<27/20$, \eqref{eq:symbol} gives
\[
 D_M-R>\frac{112}{27}-\frac{11}{10}-\frac2{25}-\frac7{32768}-\frac{27}{20}>\frac85.
\]
Every summand of $d_\Gamma(k)$ increases with $k$, proving the floor for all $n\ge M$. At $N=16384$, the same argument with $15\log2-\log3>280/27-11/10$ gives $D_n-R>39/5$ for all $n\ge N$. No source function has been numerically evaluated in these arguments.

The exact row entries, if needed for \eqref{eq:zeta}, are
\begin{align}
 Q_{nm}={}&D_n\delta_{nm}+\mathcal B_{nm}+4c_nc_m
 -\sum_{q\in\{2,3,4\}}a_qb_nb_m
   \int_{\log q-B}^B\cos(k_nx)\cos(k_m(x-\log q))\dd x,\nonumber\\
 \mathcal B_{nm}={}&b_nb_m(-1)^{n+m}\sum_{j\ge0}
 \frac{\alpha_j^2(1-e^{-2\alpha_jB})}
      {(\alpha_j^2+k_n^2)(\alpha_j^2+k_m^2)},\label{eq:rows}\\
 c_n={}&\frac{b_n(-1)^n\sinh(B/2)}{2(1/4+k_n^2)}.\nonumber
\end{align}
The lower endpoint for $q=2$ is negative and must remain signed. Formula~\eqref{eq:rows} is an exact expression, not a finite truncation prescription.

Finally, zero-extended low cosines belong to the minimum domain. Cutting them off in endpoint strips of width $\epsilon$ gives squared $L^2$ error $O(\epsilon)$ and Gamma energy $O(\epsilon(1+|\log\epsilon|))$: use a derivative bound for shifts below $\epsilon$, a support bound for larger shifts, and $\rho(v)=O(1/v)$ near zero. Their Gamma operator images have at most logarithmic endpoint singularities, hence lie in $L^2$; the remaining operator terms are bounded. Thus each cross column lies in $H$, $C_0$ converges, finite low subtraction preserves the minimum domain, and the form Schur split used above is legitimate. Indicators and discontinuous weights were used only in bounded $L^2$ statements.

\clearpage
\section*{Research status and reproducibility}
This note contains analytic proofs and a small exact-rational arithmetic script. Its reproduction bundle includes the separate frozen failed implementation and a path-redacted error record, clearly separated from the analytic proof. Neither the manuscript build nor the scalar-check command runs that implementation. A separate fixed-parameter implementation with $B=3/4$, $M=32$, $N=16384$, $J=32$ and 256-bit arithmetic did not pass its LDL positivity gate. It therefore produced no successful local-positivity certificate. No certified nonvanishing or quantitative size of the selected $\zeta$ is reported here. Failure of this sufficient comparison does not exhibit a negative direction of the original form and is not a counterexample to RH.

The implementation provenance is the SHA256 digest obtained by concatenating the following two lines:
\begin{center}\small\ttfamily
8348b5cbb0f88d5c9f34f3553ef4730c\\
625faa7d8ff5f9d36d6db3ad31298acb
\end{center}
This outcome is recorded for transparency; the scalar-check script does not reproduce or certify that separate numerical run. The failed implementation retains its original preparatory header to preserve the source digest; the accompanying status file supersedes that historical header. The analytic inequalities proved here are unaffected by its failure to close the finite positivity gate.

AI tools assisted with derivation, drafting and internal checking; responsibility for the claims and their stated limits remains with the author. Independent human review and formal proof verification have not been obtained.

\begin{thebibliography}{9}
\bibitem{FS} R. L. Frank and R. Seiringer, \emph{Non-linear ground state representations and sharp Hardy inequalities}, 2008. \href{https://arxiv.org/abs/0803.0503}{arXiv:0803.0503}.
\bibitem{DSS} G. Dusson, I. M. Sigal and B. Stamm, \emph{The Feshbach--Schur map and perturbation theory}, 2021. \href{https://arxiv.org/abs/2105.02058}{arXiv:2105.02058}.
\end{thebibliography}
\end{document}
